% 乔治·伽莫夫（综述）
% license CCBYSA3
% type Wiki

本文根据 CC-BY-SA 协议转载翻译自维基百科\href{https://en.wikipedia.org/wiki/George_Gamow}{相关文章}

乔治·伽莫夫（George Gamow，有时写作 Gammoff；原名 Гео́ргий Анто́нович Га́мов，Georgiy Antonovich Gamov，1904 年 3 月 4 日－1968 年 8 月 19 日）是一位苏联和美国的博学家、理论物理学家与宇宙学家。他是乔治·勒梅特大爆炸理论的早期倡导者与发展者之一。伽莫夫首先通过量子隧穿理论解释了α衰变的机理，提出了液滴模型——第一个原子核的数学模型，并研究了放射性衰变、恒星形成、恒星核合成、以及大爆炸核合成（他统称为“核宇宙成因论”，nucleocosmogenesis），并预言了宇宙微波背景辐射的存在以及分子遗传学的基础。伽莫夫是量子隧穿现象发展与理解的关键人物之一。

在中晚期生涯中，伽莫夫将大量精力投入教学，并撰写了多部科普著作，包括《一二三……无穷大》以及“汤普金斯先生”系列丛书（1939–1967）。他的一些著作在最初出版半个多世纪后仍在印行。为纪念他，美国科罗拉多大学博尔德分校设立了乔治·伽莫夫纪念讲座。
\subsection{早年生平与职业生涯}
伽莫夫出生于俄国帝国敖德萨（今乌克兰敖德萨）。他的父亲在中学教授俄语与文学，母亲则在女子学校教授地理与历史。除了俄语外，伽莫夫还从母亲那里学会了一些法语，并由家庭教师学习了德语。在大学期间，他又学习了英语，并最终能流利使用。伽莫夫早期的大部分论文以德语或俄语撰写，后来则以英语发表技术论文和科普著作。

他先后就读于敖德萨物理与数学学院（1922–1923），以及列宁格勒大学（1923–1929）。在列宁格勒期间，伽莫夫师从亚历山大·弗里德曼，直到弗里德曼于1925年早逝，被迫更换导师。在大学时期，伽莫夫结识了三位理论物理专业的同学——列夫·朗道、德米特里·伊万年科和马特维·布龙施泰因。四人组成了一个自称为“三剑客”的小组，定期聚会讨论并分析当时发表的开创性量子力学论文。后来，伽莫夫也用同样的称呼来形容他与拉尔夫·阿尔费和罗伯特·赫尔曼共同组成的研究小组。

毕业后，伽莫夫在哥廷根从事量子理论研究，他对原子核结构的研究成果成为其博士论文的基础。此后，他于1928年至1931年在哥本哈根大学理论物理研究所工作，中途曾前往剑桥大学卡文迪许实验室，与著名物理学家欧内斯特·卢瑟福合作。在此期间，伽莫夫继续研究原子核（提出了著名的“液滴模型”），同时还与罗伯特·阿特金森和弗里茨·豪特曼斯合作，研究恒星物理学。

1931年，年仅28岁的伽莫夫被选为苏联科学院通讯院士——成为当时最年轻的院士之一。[5][a][b] 在1931年至1933年间，他在列宁格勒镭研究所物理部工作，该部门由维塔利·赫洛平领导。在伽莫夫、伊戈尔·库尔恰托夫与列夫·米索夫斯基的指导与直接参与下，欧洲第一台回旋加速器的设计工作得以展开。1932年，伽莫夫与米索夫斯基向镭研究所学术委员会提交了加速器设计草案，并获得批准。然而，这台回旋加速器直到1937年才最终建成。[6]
\subsection{放射性衰变}
20世纪早期，科学家已经知道放射性物质具有特征性的指数衰减速率，即所谓的半衰期。与此同时，人们也发现放射性辐射具有特定的能量谱。到1928年时，伽莫夫在哥廷根成功用量子隧穿理论解释了原子核的α衰变现象，并在数学家尼古拉·科钦（的协助下完成了严格的推导。[7][8] 这一问题同时也被罗纳德·W·格尼与爱德华·U·康登独立解决。[9][10] 然而，格尼与康登的结果未能达到伽莫夫那样的定量精度。

在经典物理学中，粒子被束缚在原子核内，是因为要逃离强大的核势阱所需能量极高。按经典观点，必须施加巨大的外力才能将原子核分离，因此这种过程不可能自发发生。
然而，在量子力学框架下，粒子存在一定的概率可以“穿透”势阱的势垒并逃逸。伽莫夫在理论上求解了一个原子核的模型势场，并从第一性原理出发推导出$\alpha$衰变的半衰期与放射能量**之间的定量关系，这一关系此前是通过实验经验总结得到的，即著名的盖革–纳塔尔定律。[11]

若干年后，学界将粒子穿透库仑势垒并发生核反应的概率称为伽莫夫因子或伽莫夫–索末菲因子，以纪念他在量子隧穿理论发展中的奠基性贡献。
\subsection{叛逃}
\begin{figure}[ht]
\centering
\includegraphics[width=8cm]{./figures/f7dbe0721f3a0848.png}
\caption{1931年威廉·布拉格实验室工作人员合影：坐于中央者为威廉·亨利·布拉格，最左为物理学家A. 列别杰夫，最右为乔治·伽莫夫。} \label{fig_QZjmf_1}
\end{figure}
伽莫夫在决定逃离苏联之前，曾在多个苏联科研机构任职。然而，随着政治压迫的日益加剧，他最终选择出逃。1931年，他被正式拒绝前往意大利参加一次科学会议。同年，他与苏联物理学家**柳博芙·沃赫明采娃结婚——伽莫夫亲昵地称她为“Rho”，取自希腊字母 ρ 的发音。此后近两年间，伽莫夫和妻子不断尝试离开苏联，无论是否得到官方许可。在这段时间里，尼尔斯·玻尔及其他朋友多次邀请伽莫夫访问他们所在的国家，但他始终未能获得出境许可。

伽莫夫后来回忆说，他与妻子的前两次叛逃尝试发生在1932年，两次都打算划皮艇出逃：第一次计划横渡黑海前往土耳其，行程约250公里；第二次则打算从摩尔曼斯克出发前往挪威。然而，两次计划都因恶劣天气而失败——幸运的是，他们的举动并未被当局察觉。[12]

1933年，伽莫夫突然被批准前往比利时布鲁塞尔参加第七届索尔维会议。他坚决要求妻子必须同行，甚至表示如果不能与妻子一同前往，他将拒绝单独出席。最终，苏联当局让步，为他们夫妇二人签发了护照。

伽莫夫夫妇顺利参加了会议，并在玛丽·居里及其他物理学家的帮助下，设法延长了在西欧的停留时间。在接下来的一年里，伽莫夫先后在居里研究所、伦敦大学以及密歇根大学获得了临时研究职位。
\subsection{大爆炸核合成}
\begin{figure}[ht]
\centering
\includegraphics[width=6cm]{./figures/4e40d7f912941b32.png}
\caption{} \label{fig_QZjmf_2}
\end{figure}
伽莫夫的研究推动了炽热大爆炸理论对宇宙膨胀起源的形成与发展。他是最早运用亚历山大·弗里德曼和乔治·勒梅特提出的非静态爱因斯坦引力方程解的科学家之一——这些解描述了一个物质密度均匀、空间曲率恒定的宇宙。伽莫夫的重要突破在于，他为勒梅特关于“原初量子”的思想赋予了物理实在性：他假设在宇宙早期阶段，辐射的能量密度远高于物质的能量密度。[17]

伽莫夫的理论奠定了后世宇宙学研究的基础。他将模型进一步应用于化学元素的起源[18]以及物质凝聚形成星系的问题。[19] 在此过程中，他能够用基本物理常数——如光速 c、万有引力常数 G、精细结构常数$\alpha$以及普朗克常数 h——计算出星系的质量与尺度。

伽莫夫对宇宙学的兴趣源自他早期对恒星内部能量生成机制与元素的形成与转化的研究。[20][21][22] 而这一研究方向又起源于他最初的发现：量子隧穿效应是核$\alpha$衰变的机制。伽莫夫随后将该理论应用于相反的过程——即热核反应速率的计算，从而建立了从微观核物理到宏观宇宙演化的深刻联系。

起初，伽莫夫认为，在宇宙早期那种极高温与高密度的环境中，所有化学元素都有可能被合成出来。然而，后来他根据弗雷德·霍伊尔等人提出的强有力证据修正了这一观点——即比锂更重的元素主要是在恒星内部的热核反应以及超新星爆发过程中产生的。

伽莫夫为他提出的元素合成过程建立了一组耦合微分方程，并将这些方程的数值求解作为博士生拉尔夫·阿尔费的论文题目。伽莫夫与阿尔费的研究成果于1948年发表，论文署名为阿尔费–贝特–伽莫夫，幽默地让三位作者的首字母拼成“$\alpha$–$\beta$–$
\gamma$”，象征希腊字母顺序。[23]

在兴趣转向遗传密码问题之前，伽莫夫共发表了约二十篇宇宙学论文。最早的一篇发表于1939年，与爱德华·泰勒合作，探讨星系的形成问题。[24] 随后在1946年发表了关于宇宙核合成的首篇论文。他还撰写了多篇科普文章及学术教材，涵盖宇宙学及其他物理主题。[25]

1948年，他发表了一篇论文，研究经简化后的那组耦合方程，用以描述热中子生成质子与氘核的过程。通过合理化简并结合观测到的氢与重元素的丰度比，他得以推算出核合成开始时的物质密度，并据此计算出早期星系的质量与尺度。[26]1953年，他基于另一种方式计算出物质与辐射密度相等时的条件，得出了类似的结果。[27] 在这篇论文中，伽莫夫进一步估算了宇宙遗迹背景辐射的密度，并由此预测其当前温度约为7开尔文——这一数值约为现代实测值的两倍多。

1967年，伽莫夫发表了一篇回顾性文章，总结了他本人以及阿尔费与罗伯特·赫尔曼的相关研究成果。[28] 促使他撰写此文的契机，是彭齐亚斯与威尔逊于1965年发现了宇宙微波背景辐射。伽莫夫、阿尔费与赫尔曼认为，他们在理论上早已预言了这种辐射的存在及其来源，却未能得到应有的认可。伽莫夫尤其感到失望的是，后来发表的解释彭齐亚斯和威尔逊观测意义的论文[29]，并未引用他与合作者的早期工作，仿佛他们的理论贡献被历史遗忘。
\subsection{DNA 与 RNA}
1953年，弗朗西斯·克里克、詹姆斯·沃森、莫里斯·威尔金斯与罗莎琳·富兰克林共同发现了 DNA 大分子的双螺旋结构。此后，伽莫夫开始尝试解决一个关键问题——DNA 链中由四种碱基（腺嘌呤A、胞嘧啶C、胸腺嘧啶T、鸟嘌呤G）的排列顺序，究竟是如何控制蛋白质由氨基酸合成的。[30] 克里克曾表示，伽莫夫的设想“启发了他对遗传密码问题的思考”。[31]

据克里克回忆，[32] 伽莫夫注意到，若将四种碱基每三种一组进行组合，共有 4³ = 64 种排列方式；但如果不考虑顺序，则仅剩下 20 种独特组合。[33] 伽莫夫据此提出，这 20 种组合可能正好对应生物体中存在的**二十种氨基酸**，而所有蛋白质或许都是由这二十种氨基酸构成的。伽莫夫关于**遗传编码问题**的研究，促生了后来的**生物“简并性”模型（biological degeneracy）**。[34][35]

不过，伽莫夫提出的具体体系（被称为“**伽莫夫钻石（Gamow’s diamonds）**”）后来被证明是错误的。他假设遗传密码中的**三联体是重叠的**——例如在序列 GGAC 中，GGA 可代表一种氨基酸，而 GAC 代表另一种；同时，他还认为遗传密码**不具简并性**，即每个氨基酸只对应唯一的三碱基组合（顺序无关）。然而，后来的**蛋白质测序研究**证明这种模型并不成立。真实的遗传密码是**非重叠的**且**具简并性**——改变碱基顺序确实会改变生成的氨基酸。

1954年，伽莫夫与沃森共同创立了一个名为“**RNA 领带俱乐部（RNA Tie Club）**”的学术讨论组织，成员多为研究**遗传密码问题**的顶尖科学家，其中包括物理学家**爱德华·泰勒（Edward Teller）**和**理查德·费曼（Richard Feynman）**。沃森后来在自传中承认，伽莫夫这一**富有远见的倡议**具有重要意义。然而，这并未阻止他以半调侃的口吻形容伽莫夫是一个“**古怪离奇的巨人**”——一个喜欢变魔术、吟打油诗、豪饮作乐、恶作剧不断的天才。[36]
